\section{Limitations and Future Work}

Our experiments show a substantial difference in the performance of
our repair algorithm depending on the specific algorithms we
chose. Given the myriad classification algorithms used in practice,
there is a clear need for a future systematic study of the
relationship between dataset features, algorithms, and repair
performance.

%% \csnote{We only do ``specific algorithms''. Should we try eg. random forests?}

In addition, our discussion of disparate impact is necessarily tied to
the legal framework as defined in United States law. It would be
valuable in future work to collect the legal frameworks of different
jurisdictions, and investigate whether a single unifying formulation
is possible.

%% \csnote{Legal frameworks in Europe?}

Finally, we note that the algorithm we present operates only on numerical
attributes. Although we are satisfied with its performance, we chose
this setting mostly for its relative theoretical simplicity. A natural
avenue for future work is to investigate generalizations of our repair
procedures for datasets with different attribute types, such as
categorical data, vector-valued attributes, etc.

%% \csnote{Numeric data vs categorical attributes. The authors need to spend more time discussing: "Specifically, our goal will be to preserve the relative per-attribute ordering as follows." Why is that the goal? What other potential goals could there have been?}
